% GENERATED FILE. Edit canonical lessons, then run npm run generate:books.
% canonical-source-hash: fnv1a64:36eaf702c7d91cc1
% canonical-lessons: PA-C249-balna-kindle, PA-C249-phukna, PA-C249-dhumam, PA-C249-kola, PA-C249-angiar

\chapter{Fire and Smoke}
\label{ch:pa-a2-fire-and-smoke}
\hlchaptermodality{249}

\begin{chapteropening}
\textbf{By the end of this chapter:} \emph{I can talk about lighting a kitchen fire.}

Complete the last lesson of chapter 249: I can talk about lighting a kitchen fire.
\end{chapteropening}

\section[\texorpdfstring{\pa{ਬਾਲਣਾ} (bālṇā)}{ਬਾਲਣਾ (bālṇā)}]{\pa{ਬਾਲਣਾ} (bālṇā) — to light a fire; to light a lamp}

\label{lesson:PA-C249-balna-kindle}



\begin{quote}
Before the new one: say the Punjabi for a chain, then the Punjabi for a wall peg.
\end{quote}

\subsection*{You'll want to know: \pa{ਬਾਲਣਾ}}
\textbf{\pa{ਬਾਲਣਾ}} — \emph{bālṇā} — \textquotedblleft{}to light a fire; to light a lamp\textquotedblright{}.

\begin{quote}
\pa{ਦਾਦੀ} \pa{ਸਵੇਰੇ} \pa{ਚੁੱਲ੍ਹਾ} \pa{ਬਾਲਦੀ} \pa{ਹੈ।}
\end{quote}

\begin{quote}
\emph{dādī savere chullhā bāldī hai.}
\end{quote}

\begin{quote}
Grandmother lights the stove in the morning.
\end{quote}

Past (I …): \pa{ਮੈਂ} \pa{ਦੀਵਾ} \pa{ਬਾਲਿਆ}, \emph{maiṁ dīvā bāliā}.

\begin{cousinweb}
\textbf{\pa{ਬਾਲਣਾ}} (\emph{bālṇā}) is inherited from a reconstructed \emph{*dvālayati}, \textquotedblleft{}kindles\textquotedblright{}; CDIAL lists Punjabi \emph{bālaṇā}, \textquotedblleft{}to kindle\textquotedblright{} (CDIAL 6671).
\end{cousinweb}

\begin{grammarlens}[title={What you've built}]
One more word for this chapter.
\end{grammarlens}

\subsection*{Your turn}
\begin{itemize}
\raggedright
  \item \emph{Say it:} \emph{bālṇā}
  \item \emph{Say it:} \emph{bālṇā}, once more
  \item \emph{Say it:} \emph{dādī savere chullhā bāldī hai.} — Grandmother lights the stove in the morning.
  \item \emph{Recall:} say the Punjabi for a chain, then the Punjabi for a wall peg, then say \emph{bālṇā} again
\end{itemize}

\begin{tcolorbox}[breakable,colback=teal!4,colframe=teal!35!black,title={Before you move on}]
What does \pa{ਬਾਲਣਾ} mean? (\textquotedblleft{}To light a fire; to light a lamp\textquotedblright{}.) Say it once more.
\end{tcolorbox}
\section[\texorpdfstring{\pa{ਫੂਕਣਾ} (phūkṇā)}{ਫੂਕਣਾ (phūkṇā)}]{\pa{ਫੂਕਣਾ} (phūkṇā) — to blow on a fire}

\label{lesson:PA-C249-phukna}



\begin{quote}
Before the new one: say the Punjabi for a wall peg, then the Punjabi for to light a fire.
\end{quote}

\subsection*{You'll want to know: \pa{ਫੂਕਣਾ}}
\textbf{\pa{ਫੂਕਣਾ}} — \emph{phūkṇā} — \textquotedblleft{}to blow on a fire\textquotedblright{}.

\begin{quote}
\pa{ਸਵੇਰੇ} \pa{ਮੈਂ} \pa{ਚੁੱਲ੍ਹਾ} \pa{ਫੂਕਦੀ} \pa{ਹਾਂ।}
\end{quote}

\begin{quote}
\emph{savere maiṁ chullhā phūkdī hāṁ.}
\end{quote}

\begin{quote}
In the morning I blow on the stove to get the fire going.
\end{quote}

Past (I …): \pa{ਮੈਂ} \pa{ਚੁੱਲ੍ਹਾ} \pa{ਫੂਕਿਆ}, \emph{maiṁ chullhā phūkiā}.

\begin{cousinweb}
\textbf{\pa{ਫੂਕਣਾ}} (\emph{phūkṇā}) is inherited from a reconstructed \emph{*phūtka}, \textquotedblleft{}blowing\textquotedblright{}; CDIAL lists Punjabi \emph{phūkṇā}, \textquotedblleft{}to blow\textquotedblright{} (CDIAL 9102).
\end{cousinweb}

\begin{grammarlens}[title={What you've built}]
One more word for this chapter.
\end{grammarlens}

\subsection*{Your turn}
\begin{itemize}
\raggedright
  \item \emph{Say it:} \emph{phūkṇā}
  \item \emph{Say it:} \emph{phūkṇā}, once more
  \item \emph{Say it:} \emph{savere maiṁ chullhā phūkdī hāṁ.} — In the morning I blow on the stove to get the fire going.
  \item \emph{Recall:} say the Punjabi for a wall peg, then the Punjabi for to light a fire, then say \emph{phūkṇā} again
  \item \emph{Write it:} copy \textbf{\pa{ਫੂਕਣਾ}} once, checking the whole word
\end{itemize}

\begin{tcolorbox}[breakable,colback=teal!4,colframe=teal!35!black,title={Before you move on}]
What does \pa{ਫੂਕਣਾ} mean? (\textquotedblleft{}To blow on a fire\textquotedblright{}.) Say it once more.
\end{tcolorbox}
\section[\texorpdfstring{\pa{ਧੂੰਆਂ} (dhūṁāṁ)}{ਧੂੰਆਂ (dhūṁāṁ)}]{\pa{ਧੂੰਆਂ} (dhūṁāṁ) — smoke}

\label{lesson:PA-C249-dhumam}



\begin{quote}
Before the new one: say the Punjabi for to light a fire, then the Punjabi for to blow on a fire.
\end{quote}

\subsection*{You'll want to know: \pa{ਧੂੰਆਂ}}
\textbf{\pa{ਧੂੰਆਂ}} — \emph{dhūṁāṁ} — \textquotedblleft{}smoke\textquotedblright{}.

\begin{quote}
\pa{ਗਿੱਲੀ} \pa{ਲੱਕੜ} \pa{ਨਾਲ} \pa{ਬਹੁਤ} \pa{ਧੂੰਆਂ} \pa{ਹੁੰਦਾ} \pa{ਹੈ।}
\end{quote}

\begin{quote}
\emph{gillī lakkaṛ nāl bahut dhūṁāṁ hundā hai.}
\end{quote}

\begin{quote}
Wet wood makes a lot of smoke.
\end{quote}

A masculine noun.

\begin{cousinweb}
\textbf{\pa{ਧੂੰਆਂ}} (\emph{dhūṁāṁ}) is inherited from Sanskrit \emph{dhūma}, \textquotedblleft{}smoke, vapour\textquotedblright{}; CDIAL lists Punjabi \emph{dhūā̃}, \textquotedblleft{}smoke\textquotedblright{} (CDIAL 6849).
\end{cousinweb}

\begin{grammarlens}[title={What you've built}]
One more word for this chapter.
\end{grammarlens}

\subsection*{Your turn}
\begin{itemize}
\raggedright
  \item \emph{Say it:} \emph{dhūṁāṁ}
  \item \emph{Say it:} \emph{dhūṁāṁ}, once more
  \item \emph{Say it:} \emph{gillī lakkaṛ nāl bahut dhūṁāṁ hundā hai.} — Wet wood makes a lot of smoke.
  \item \emph{Recall:} say the Punjabi for to light a fire, then the Punjabi for to blow on a fire, then say \emph{dhūṁāṁ} again
  \item \emph{Write it:} copy \textbf{\pa{ਧੂੰਆਂ}} once, checking the whole word
\end{itemize}

\begin{tcolorbox}[breakable,colback=teal!4,colframe=teal!35!black,title={Before you move on}]
What does \pa{ਧੂੰਆਂ} mean? (\textquotedblleft{}Smoke\textquotedblright{}.) Say it once more.
\end{tcolorbox}
\section[\texorpdfstring{\pa{ਕੋਲਾ} (kolā)}{ਕੋਲਾ (kolā)}]{\pa{ਕੋਲਾ} (kolā) — charcoal; coal}

\label{lesson:PA-C249-kola}



\begin{quote}
Before the new one: say the Punjabi for to blow on a fire, then the Punjabi for smoke.
\end{quote}

\subsection*{You'll want to know: \pa{ਕੋਲਾ}}
\textbf{\pa{ਕੋਲਾ}} — \emph{kolā} — \textquotedblleft{}charcoal; coal\textquotedblright{}.

\begin{quote}
\pa{ਕੋਲਾ} \pa{ਕਾਲਾ} \pa{ਅਤੇ} \pa{ਹੌਲਾ} \pa{ਹੈ।}
\end{quote}

\begin{quote}
\emph{kolā kālā ate haulā hai.}
\end{quote}

\begin{quote}
Charcoal is black and light.
\end{quote}

A masculine noun.

\begin{cousinweb}
\textbf{\pa{ਕੋਲਾ}} (\emph{kolā}) is inherited: CDIAL lists Punjabi \emph{kolā}, \textquotedblleft{}dead coal, charcoal\textquotedblright{}, under a late Sanskrit \emph{kōkila}, \textquotedblleft{}lighted coal\textquotedblright{}; the entry records the view that this Sanskrit word was taken from Middle Indo-Aryan and, before that, perhaps from a Munda language (CDIAL 3484).
\end{cousinweb}

\begin{grammarlens}[title={What you've built}]
One more word for this chapter.
\end{grammarlens}

\subsection*{Your turn}
\begin{itemize}
\raggedright
  \item \emph{Say it:} \emph{kolā}
  \item \emph{Say it:} \emph{kolā}, once more
  \item \emph{Say it:} \emph{kolā kālā ate haulā hai.} — Charcoal is black and light.
  \item \emph{Recall:} say the Punjabi for to blow on a fire, then the Punjabi for smoke, then say \emph{kolā} again
  \item \emph{Write it:} copy \textbf{\pa{ਕੋਲਾ}} once, checking the whole word
\end{itemize}

\begin{tcolorbox}[breakable,colback=teal!4,colframe=teal!35!black,title={Before you move on}]
What does \pa{ਕੋਲਾ} mean? (\textquotedblleft{}Charcoal; coal\textquotedblright{}.) Say it once more.
\end{tcolorbox}
\section[\texorpdfstring{\pa{ਅੰਗਿਆਰ} (aṅgiār)}{ਅੰਗਿਆਰ (aṅgiār)}]{\pa{ਅੰਗਿਆਰ} (aṅgiār) — a glowing coal; an ember}

\label{lesson:PA-C249-angiar}



\begin{quote}
Before the new one: say the Punjabi for smoke, then the Punjabi for charcoal.
\end{quote}

\subsection*{You'll want to know: \pa{ਅੰਗਿਆਰ}}
\textbf{\pa{ਅੰਗਿਆਰ}} — \emph{aṅgiār} — \textquotedblleft{}a glowing coal; an ember\textquotedblright{}.

\begin{quote}
\pa{ਚੁੱਲ੍ਹੇ} \pa{ਵਿੱਚ} \pa{ਲਾਲ} \pa{ਅੰਗਿਆਰ} \pa{ਹਨ।}
\end{quote}

\begin{quote}
\emph{chullhe vich lāl aṅgiār han.}
\end{quote}

\begin{quote}
There are red embers in the stove.
\end{quote}

A masculine noun.

\begin{cousinweb}
\textbf{\pa{ਅੰਗਿਆਰ}} (\emph{aṅgiār}) is inherited from Sanskrit \emph{aṅgāra}, \textquotedblleft{}glowing charcoal\textquotedblright{}; CDIAL lists Punjabi \emph{aṅgyār}, \textquotedblleft{}charcoal\textquotedblright{} (CDIAL 125).
\end{cousinweb}

\begin{grammarlens}[title={What you've built}]
One more word for this chapter.
\end{grammarlens}

\subsection*{Your turn}
\begin{itemize}
\raggedright
  \item \emph{Say it:} \emph{aṅgiār}
  \item \emph{Say it:} \emph{aṅgiār}, once more
  \item \emph{Say it:} \emph{chullhe vich lāl aṅgiār han.} — There are red embers in the stove.
  \item \emph{Recall:} say the Punjabi for smoke, then the Punjabi for charcoal, then say \emph{aṅgiār} again
  \item \emph{Write it:} copy \textbf{\pa{ਅੰਗਿਆਰ}} once, checking the whole word
\end{itemize}

\begin{tcolorbox}[breakable,colback=teal!4,colframe=teal!35!black,title={Before you move on}]
What does \pa{ਅੰਗਿਆਰ} mean? (\textquotedblleft{}A glowing coal; an ember\textquotedblright{}.) Say it once more.
\end{tcolorbox}
